\begin{textsample}{NEG Dim 2 – global\_north\_2025\_01 – Score 1.00 – global\_north\_2025\_01/120719620.txt}  \label{ex:f2_neg_048}
This is the last article you can read this month You can read more article this month You can read more articles this month Sorry your limit is up for this month THIS year Welcome to the Fringe Palestine will bring 16 Palestinian artists to perform at the Edinburgh Fringe, to platform their voices and celebrate Palestinian art and culture, with freedom and without censorship.
It will be a chance for poets, theatre-makers, comedians and dancers from across Palestine and the diaspora to engage with the festival, meet audiences, see shows, make friends, enjoy the city and share their work at the biggest arts festival on Earth.
Welcome to the Fringe is crowdfunding to make this happen.
As an independent group of Scotland-based artists and producers who have worked regularly in Palestine and the wider Middle East over many decades, Welcome to the Fringe wants to embrace the lessons in hospitality, culture and liberation that they have learnt from Palestinians and bring that spirit to the Festival.
Welcome @ @ @ @ @ @ @ @ @ @ 2025.
There will be art and politics, as well as music, food, drink, stories, jokes and conversation.
They need donations, volunteers, and accommodation to make it happen.
The Edinburgh International Festival was founded in 1947 to"provide a platform for the flowering of the human spirit"and bring the people of Europe together following the second world war.
In 1951 Hamish Henderson, Joe Corrie, Janey Buchan and others organised The Edinburgh People ' s Festival -- funded by the Communist Party and the Scottish Trade Union Congress -- this was a celebration of working-class art and politics and an antidote to the elitism of the established Festival.
A truly internationalist ceilidh where folk music, poetry and comedy flourished : the People ' s Festival was the forerunner of today ' s Edinburgh Fringe.
Welcome to the Fringe deeply values the legacy and aims of these two festivals, and celebrates the Fringe as a place of welcome, refuge, and celebration of the human spirit today.
It is essential @ @ @ @ @ @ @ @ @ @ apartheid and ethnic cleansing, attempts to dehumanise Palestinians have reached fever pitch across the last 15 months of Israel ' s genocidal war in Gaza.
Now more than ever we need to hear from Palestinians, and to take inspiration from their work.
This is a small stand against the brutal suppression and destruction faced by Palestinian artists at the hands of the Israeli state, and the ongoing censorship, restrictions and barriers placed in the way of Palestinian art and artists by the UK and the EU.
It is a chance to do something, however small it seems, to combat the ongoing dehumanisation and violence against the Palestinian people.
Above all, Welcome To The Fringe Palestine aims to listen and learn from Palestinians.
Currently it is not possible for artists from Gaza to travel to the UK, with the British Home Office offering no access to visas.
Similarly Gazans taking refuge in Egypt are being denied paperwork to travel.
Artists in the West Bank are having visas denied and are being blocked by Israel @ @ @ @ @ @ @ @ @ @ and ensuring that artists from Palestine are not denied access to the Edinburgh Festival and that Edinburgh audiences are not denied access to Palestinian art is an essential aim.
Palestinian dancer and co-organiser of Welcome to the Fringe Farah Saleh says :"It is possible that many Palestinians will simply have no opportunity to leave Palestine to come to Edinburgh.
Visa restrictions by UK and Israeli authorities could leave artists unable to get to Edinburgh.
Increasing Israeli violence makes travel more difficult.
In those circumstances we will need to explore ways of funding Palestinian voices to be heard in Edinburgh via other means and to raise the question of who the Edinburgh festival is for."I am proud to be part of the committee organising Welcome to the Fringe.
Every one of us is working voluntarily and we are fundraising from scratch.
To make the showcase happen we are reliant on the generosity, skills and solidarity of the many.
And so, if you are interested in helping with the Welcome To The Fringe Palestine project visit @ @ @ @ @ @ @ @ @ @

% matched lemmas: 
\end{textsample}
